\documentclass{article}

\usepackage{amsmath}
\usepackage{amsthm}
\usepackage{amssymb}

\usepackage[utf8]{inputenc}     
\usepackage[T1]{fontenc}
\usepackage[magyar]{babel}

\title{Alapgeometria}
\author{Rovenszky Robin\thanks{Rovenszky Robin fordításával}}
\date{Legutóbb frissítve: 2026. Március 16.}

\newtheorem{theorem}{Tétel}

\begin{document}

\maketitle
\tableofcontents
\newpage

\begin{theorem}[Pitagorasz-tétel]
Egy derékszögű háromszögben, amelynek befogói $a$ és $b$, átfogója pedig $c$, a következő összefüggés áll fenn:
\[
a^2 + b^2 = c^2.
\]
\end{theorem}

\begin{proof}
Vegyünk egy $(a+b)$ oldalhosszúságú négyzetet. Helyezzünk el a négyzet belsejében négy azonos derékszögű háromszöget, amelyek befogói $a$ és $b$, átfogójuk $c$, úgy, hogy az átfogók egy kisebb négyzetet alkossanak a középpontban.

A nagy négyzet területe:
\[
(a+b)^2.
\]

Ez a terület egyenlő a négy háromszög területének és a középső kis négyzet területének összegével.

Minden háromszög területe:
\[
\frac{1}{2}ab,
\]
így a négy háromszög együtt:
\[
4 \cdot \frac{1}{2}ab = 2ab.
\]

A középső négyzet oldala $c$, területe:
\[
c^2.
\]

Tehát
\[
(a+b)^2 = 2ab + c^2.
\]

A bal oldalt kibontva:
\[
a^2 + 2ab + b^2 = 2ab + c^2.
\]

A két oldalból kivonva $2ab$-t:
\[
a^2 + b^2 = c^2.
\]

Ez bizonyítja a tételt.
\end{proof}

\newpage

\section{Kapcsolódó feladatok}

\subsection{1. feladat}

Első Pitagorasz-tétel:
\begin{equation}
    12^2 + x^2 = 15^2
\end{equation}
\begin{equation}
    144+x^2 = 225
\end{equation}
\begin{equation}
    x^2 = 81
\end{equation}
\begin{equation}
    x = 9
\end{equation}
Második Pitagorasz-tétel:
\begin{equation}
    12^2 + y^2 = 13^2
\end{equation}
\begin{equation}
    144 + y^2 = 169
\end{equation}
\begin{equation}
    y^2 = 25
\end{equation}
\begin{equation}
    y = 5
\end{equation}
Tehát a végső eredményünk: $x + y = 14$. \qed

\subsection{2. feladat}

Első Pitagorasz-tétel:
\begin{equation}
    x^2 + (12-r)^2 = (12+r)^2
\end{equation}
Második Pitagorasz-tétel:
\begin{equation}
    x^2 + r^2 = (24-r)^2
\end{equation}
(9) - (10) értékelése:
\begin{equation}
    (12-r)^2 - r^2 = (12+r)^2 - (24-r)^2
\end{equation}
\begin{equation}
    (12-2r) \cdot 12 = 36 \cdot (2r - 12)
\end{equation}
12-vel osztva:
\begin{equation}
    (12-2r) = (2r - 12) \cdot 3
\end{equation}
Majd 2-vel osztva:
\begin{equation}
    6-r = 3 \cdot (r-6)
\end{equation}
\begin{equation}
    6-r = 3r - 18
\end{equation}
18+r hozzáadásával:
\begin{equation}
    24 = 4r
\end{equation}
\begin{equation}
    r = 6
\end{equation}
\qed

\subsection{3. feladat}
\begin{theorem}[1. tétel]
Egy paralelogrammára:
\begin{equation}
    e^2 + f^2 = a^2 \cdot2 + b^2 \cdot2
\end{equation}
\begin{proof}
Első Pitagorasz-tétel:
\begin{equation}
    (a-x)^2 + m^2 = e^2
\end{equation}
Második Pitagorasz-tétel:
\begin{equation}
    (a+x)^2 + m^2 = f^2
\end{equation}
(19) + (20) értékelése:
\begin{equation}
    a^2 - 2ax + x^2 + m^2 + a^2 + 2ax + x^2 + m^2 = e^2 + f^2
\end{equation}
Egyszerűsítve és a $m^2 + x^2 = b^2$ összefüggést felhasználva:
\begin{equation}
    2a^2 + 2b^2 = e^2 + f^2
\end{equation}
\end{proof}

\end{theorem}

\subsection{4. feladat}
Legyen a kis kör sugara $r$.\\
Első Pitagorasz-tétel:
\begin{equation}
    x^2 + (12 - r)^2 = (12 + r)^2
\end{equation}
Második Pitagorasz-tétel:
\begin{equation}
    r^2 + x^2 = (24-r)^2
\end{equation}
Lásd az (1.2) Kapcsolódó feladatok: 2. feladat megoldását. Ezek két egyenértékű feladat, 90 fokos elforgatással. A képek segítenek a tisztázásban.

\subsection{5. feladat}
Most tegyük a köröket téglalapok belsejébe ahelyett, hogy más körökbe tennénk!\\
Keressük a téglalap magasságát a kis körök átmérője alapján.\\
$r_0 = 12$

\subsection{6. feladat}
Egy másik kör-belső-téglalap probléma. A téglalap 60x48, mi a kis körök sugara?
\end{document}